% !TEX program = xelatex
\documentclass{misisarticle}
\grnti{28.23.37, 16.31.21}
\asjc{1702, 1706}
\oecd{1.02}
\title{Сравнительное исследование и дистилляция больших языковых моделей для классификации эмоций в русскоязычных отзывах пользователей интернет-сервисов}
\misisauthor{Верещагина Татьяна Владимировна}{БПМ-23-ПО-3, кафедра инженерной кибернетики НИТУ «МИСИС»}{Москва, Россия}{}
\abstract{В работе рассматривается определение эмоций в русскоязычных отзывах пользователей. Для решения этой задачи предлагается использовать большую языковую модель (LLM), которая размечает отзывы, и модель на основе ruBERT, которая обучается на полученной разметке. Такой подход позволяет проверить, можно ли сократить затраты на обработку текстов, сохранив качество классификации. Выделяются пять эмоций: радость, грусть, удивление, страх и гнев. Один отзыв может содержать несколько эмоций или не содержать ни одной из них. Для исследования выбраны CEDR с эмоциональной разметкой и RuReviews с отзывами, для которых такую разметку ещё нужно получить. Планируется сравнить классический метод TF-IDF с логистической регрессией, ruBERT, LLM и ruBERT после обучения на разметке LLM. В статье описаны задача, выбранные методы и порядок их проверки. Экспериментальная часть будет выполнена на следующем этапе работы.}
\keywords{анализ эмоций, русскоязычные отзывы, многометочная классификация, дистилляция знаний, большие языковые модели, ruBERT}
\begin{document}
\maketitle
\section{Введение}
Пользователи интернет-сервисов оставляют отзывы о товарах и услугах, рассказывают о проблемах и делятся впечатлениями. По этим текстам компании могут понять, что вызывает недовольство и что нравится пользователям. Но при большом количестве отзывов читать каждый из них вручную сложно, поэтому для их обработки применяют методы машинного обучения.

Обычно в отзывах определяют тональность: положительную, нейтральную или отрицательную. Однако такой оценки не всегда достаточно. Например, два отрицательных отзыва могут выражать разные эмоции: в одном пользователь злится, а в другом расстроен. В одном тексте могут сочетаться и несколько эмоций. Поэтому в этой работе рассматривается более подробная задача --- определение эмоциональных категорий отзыва.

Один из способов её решения --- большая языковая модель. Ей можно передать текст и инструкцию с описанием эмоций, не обучая отдельный классификатор. При этом обработка большого числа отзывов через LLM требует времени и вычислительных ресурсов, а при использовании платного API --- денежных затрат. Другой вариант --- использовать LLM для подготовки обучающих данных, а затем обучить на них меньшую модель. Возможность такого подхода рассматривается в работе Wang и соавторов~\cite{wang}.

Цель работы --- сравнить несколько подходов к определению эмоций в русскоязычных отзывах и проверить, насколько полезна разметка LLM для обучения ruBERT. Для этого необходимо подготовить данные, проверить качество автоматических меток, обучить модели и сравнить их на одних и тех же тестовых текстах. При сравнении будут учитываться как качество ответов, так и скорость и стоимость обработки.

\section{Анализ работ}
Для исследования эмоций в русскоязычных текстах подходит корпус CEDR, описанный Sboev и соавторами~\cite{sboev}. Он содержит предложения из LiveJournal, Lenta.ru и Twitter, размеченные по пяти эмоциям: радость, грусть, удивление, страх и гнев. Одному предложению может соответствовать несколько меток. В версии 0.1.1 содержится 9410 примеров: 7528 в обучающей части и 1882 в тестовой~\cite{cedr}. У этого корпуса уже есть нужная разметка, но его тексты отличаются от отзывов. Поэтому результат на CEDR сам по себе не покажет, насколько хорошо модель работает с отзывами пользователей.

Smetanin и Komarov исследуют анализ тональности русскоязычных товарных отзывов~\cite{smetanin}. Связанный с этой работой набор RuReviews содержит около 90 тысяч отзывов, разделённых на три сбалансированных класса тональности~\cite{rureviews}. В отличие от CEDR, в нём есть подходящий тип текстов, но нет эмоциональных меток. Положительную тональность нельзя автоматически приравнять к радости, а отрицательную --- к гневу. Для данной задачи отзывы нужно разметить заново. Кроме того, RuReviews содержит товарные отзывы одной категории. Для выводов о других интернет-сервисах потребуются дополнительные данные.

В качестве основы для классификатора выбрана модель ruBERT. Она относится к семейству BERT, в котором представление слова зависит от окружающего текста~\cite{bert,rubert}. Для определения эмоций к модели добавляется выходной слой, который оценивает каждую категорию отдельно. Благодаря этому один отзыв может получить несколько меток. ruBERT, обученный на CEDR, будет использоваться для сравнения с вариантом, которому дополнительно переданы отзывы с разметкой LLM.

В работе Измайлова и Устинова~\cite{izmailov} сравниваются две дообученные модели: rubert-tiny2 и rubert\_base\_cased\_sentence. Авторы используют объединённый корпус отзывов о фильмах и комментариев с разметкой токсичности. Более крупная модель показывает лучшее качество, тогда как компактная требует меньше времени на обучение и имеет меньший размер. Эта работа близка к выбранной теме применением ruBERT к русскоязычным пользовательским текстам и сравнением качества с вычислительными затратами. Но в ней рассматривается бинарная классификация, а в настоящем исследовании --- несколько эмоций одновременно. Поэтому её показатели нельзя напрямую использовать как ожидаемый результат на CEDR и RuReviews.

Береснев~\cite{beresnev} рассматривает предобученные модели для корпоративного чат-бота. В работе реализованы семантический поиск по документам с помощью метода ближайших соседей и BERT-подобных моделей, а также классификация текста с подбором параметров дообучения. Предлагается модульная архитектура, в которой отдельные модели отвечают за поиск, классификацию и генерацию ответа. Для настоящей работы полезен сам принцип выбора модели под конкретную задачу и доступные ресурсы. Однако разметка эмоций с помощью LLM и перенос этих меток в ruBERT в данной статье не исследуются.

В работе Wang и соавторов GPT-3 применяется для получения меток, на которых затем обучаются другие модели~\cite{wang}. Такой способ полезен, когда ручная разметка большого набора текстов занимает много времени. Однако автоматическая разметка тоже содержит ошибки, поэтому её качество нужно проверять отдельно.

Перенос знаний от большой модели к меньшей называют дистилляцией. В работе Hinton и соавторов ученик обучается по распределению ответов учителя~\cite{hinton}. В выбранном здесь варианте LLM передаёт только названия эмоций. Поэтому речь идёт об обучении на псевдометках --- категориях, которые определила другая модель, а не человек. Вероятности ответов учителя в этом варианте не используются.

Рассмотренные работы дают основу для выбора данных и методов. Теперь нужно проверить их сочетание в одной задаче: улучшит ли разметка LLM качество ruBERT на отзывах и оправдает ли возможный прирост затраты на её получение.

\section{Постановка задачи}
На вход модели поступает русскоязычный текст отзыва. На выходе требуется получить набор выраженных в нём эмоций. Задача является многометочной, поскольку категории не исключают друг друга.

Пусть $x$ --- текст отзыва, а $\mathcal{E}=\{e_1,\ldots,e_5\}$ --- пять выбранных эмоций. Тогда ответ модели можно записать в виде
\begin{equation}
f_\theta(x)=\hat{\mathbf y},\qquad \hat{\mathbf y}\in\{0,1\}^{5},
\end{equation}
где единица означает наличие соответствующей эмоции, а ноль --- её отсутствие. Если все пять значений равны нулю, текст считается не содержащим выраженных эмоций из выбранного списка. Отдельная шестая категория «нет эмоции» не вводится.

Для обучения используются два источника. Первый --- CEDR с готовыми эмоциональными метками, обозначенный $D_g$. Второй --- отзывы RuReviews, для которых LLM создаёт псевдометки. Полученный набор обозначается $D_s$; далее он также называется серебряной разметкой.

Основная задача --- сравнить модель, обученную на $D_g\cup D_s$, с ruBERT, обученным только на $D_g$, и остальными выбранными методами. Предполагается, что дополнительные отзывы помогут модели лучше работать с этим типом текстов. При этом улучшение не гарантировано: ошибки учителя могут перейти в обучающую выборку и ухудшить результат.

\section{Описание методов}
\subsection{Подготовка данных}
Тестовая часть CEDR используется только для итоговой оценки. Из обучающей части выделяется валидационная выборка, по которой подбираются параметры модели и пороги отнесения к эмоциям. Такое разделение нужно, чтобы настройка модели не зависела от ответов на тесте.

Перед разметкой RuReviews удаляются пустые отзывы, лишние пробелы, одинаковые и почти одинаковые тексты. При этом отрицания, эмодзи и знаки препинания сохраняются: они могут менять эмоциональный смысл. Например, удаление частицы «не» способно изменить смысл высказывания на противоположный.

Для ручной проверки планируется отобрать около 500 отзывов. Их независимо размечают два человека, после чего спорные случаи разбираются отдельно. Эти тексты заранее исключаются из обучающего набора с псевдометками. Часть ручной выборки используется для уточнения инструкции, а оставшаяся часть --- для окончательного тестирования. Итоговый тест не используется при выборе примеров для LLM и настройке модели. Если в нём окажется слишком мало примеров редких эмоций, выборку потребуется увеличить.

\subsection{Разметка с помощью LLM}
Большая языковая модель выступает в роли учителя. Для неё составляется инструкция с определениями эмоций и правилами выбора нескольких меток. Если ни одна эмоция не выражена, модель должна вернуть пустой список. Ответ задаётся в виде списка допустимых категорий, чтобы его можно было обработать программно.

Планируется проверить два режима. В zero-shot модель получает только инструкцию, а в few-shot к ней добавляются размеченные примеры из части данных, выделенной для разработки. Для каждого варианта сохраняются название и версия LLM, инструкция, параметры генерации и дата обращения. Это нужно, чтобы было понятно, при каких условиях получена разметка. Точная модель будет выбрана до итогового тестирования с учётом доступных ресурсов.

Качество меток LLM проверяется по ручной разметке отдельно для каждой эмоции. Для сравнения используются F1 и коэффициент согласия Коэна $\kappa$. Согласие двух людей рассчитывается до разбора спорных случаев. Также учитывается число примеров каждой категории: высокий общий показатель может скрывать ошибки на редких эмоциях. Некорректные ответы LLM и повторные запросы сохраняются, а затраты на них включаются в общую стоимость разметки.

\subsection{Обучение ruBERT}
К ruBERT добавляется выходной слой с пятью значениями --- по одному для каждой эмоции. Сигмоидальная функция преобразует их в оценки от нуля до единицы. Они рассматриваются независимо, поэтому наличие одной эмоции не мешает выбрать другую.

Для обучения используется бинарная кросс-энтропия:
\begin{equation}
\ell(\mathbf y,\mathbf p)=-\frac{1}{5}\sum_{j=1}^{5}\left[y_j\log p_j+(1-y_j)\log(1-p_j)\right],
\end{equation}
где $y_j$ --- метка эмоции, а $p_j$ --- оценка модели. Ошибка рассчитывается для каждой категории и затем усредняется.

Поскольку ручные метки и псевдометки могут отличаться по качеству, их вклад в обучение задаётся отдельно:
\begin{equation}
\mathcal L=(1-\lambda)\frac{1}{|D_g|}\sum_{D_g}\ell+\lambda\frac{1}{|D_s|}\sum_{D_s}\ell,
\quad 0\leq\lambda\leq1.
\end{equation}
Параметр $\lambda$ определяет вес серебряной разметки. Его значение и число дополнительных отзывов подбираются на валидационной выборке. По ней же выбираются пороги $\tau_j$: эмоция включается в ответ, если оценка модели не ниже соответствующего порога. Отдельно можно проверить последовательное обучение --- сначала на CEDR, затем на псевдометках. Его результаты будут учитываться как самостоятельный вариант.

\subsection{Сравнение моделей}
Для сравнения выбраны TF-IDF с логистической регрессией, ruBERT на CEDR, LLM в режимах zero-shot и few-shot и ruBERT с дополнительной разметкой LLM. В классическом методе для каждой из пяти эмоций обучается отдельный бинарный классификатор.

Основной метрикой будет macro-F1 --- среднее значение F1 по пяти категориям. Оно позволяет учитывать каждую эмоцию с одинаковым весом. Дополнительно рассчитываются micro-F1, precision, recall и F1 для отдельных эмоций, а также доля текстов, для которых полностью совпал набор меток. Если в тесте нет положительных примеров какой-либо эмоции, это указывается при представлении результатов. Оценки на CEDR и на ручной выборке отзывов приводятся отдельно, поскольку эти данные отличаются по источнику и типу текстов.

Чтобы результат обучения ruBERT не зависел от одного случайного запуска, планируется выполнить не менее трёх запусков с разными значениями seed. По ним рассчитываются среднее значение и стандартное отклонение метрик. Для каждого запуска сохраняются версия модели, токенизатор, максимальная длина текста, размер батча, число эпох и условие остановки обучения.

Помимо качества сравнивается время обработки отзывов. При замерах указываются оборудование, использование CPU или GPU, длина текстов и размер батча. Перед измерениями выполняется прогрев модели. Затем рассчитываются медиана и 95-й перцентиль задержки, число обработанных текстов в секунду и пиковое потребление памяти. Для LLM через API отдельно учитывается влияние сетевого запроса.

Стоимость сравнивается для одинакового числа отзывов. Затраты на получение псевдометок рассчитываются отдельно от затрат на дальнейшее применение ruBERT. Для платного API указываются тариф и дата расчёта. Это позволит оценить, при каком объёме обработки обучение меньшей модели может оказаться выгоднее постоянных запросов к LLM.

Дополнительно планируется проверить, как меняется результат при увеличении числа размеченных LLM отзывов и изменении веса $\lambda$. Ошибки будут разобраны на текстах с сарказмом, отрицаниями и несколькими эмоциями. По результатам экспериментов можно будет сделать вывод, помогает ли выбранный способ дистилляции в этой задаче. Пока в статье представлены подготовка исследования и методы; численные результаты ещё не получены.

\begin{thebibliography}{99}
\bibitem{sboev} Sboev A., Naumov A., Rybka R. Data-Driven Model for Emotion Detection in Russian Texts // Procedia Computer Science. 2021. Vol. 190. P. 637--642. DOI: 10.1016/j.procs.2021.06.075.
\bibitem{cedr} CEDR: Corpus for Emotions Detecting in Russian-language text sentences. Версия 0.1.1. URL: \url{https://github.com/sag111/CEDR} (дата обращения: 05.11.2026).
\bibitem{smetanin} Smetanin S., Komarov M. Sentiment Analysis of Product Reviews in Russian using Convolutional Neural Networks // 2019 IEEE 21st Conference on Business Informatics. 2019. P. 482--486. DOI: 10.1109/CBI.2019.00062.
\bibitem{rureviews} RuReviews: An Automatically Annotated Sentiment Analysis Dataset for Product Reviews in Russian. URL: \url{https://github.com/sismetanin/rureviews} (дата обращения: 05.11.2026).
\bibitem{bert} Devlin J., Chang M.-W., Lee K., Toutanova K. BERT: Pre-training of Deep Bidirectional Transformers for Language Understanding // NAACL-HLT. 2019. P. 4171--4186. DOI: 10.18653/v1/N19-1423.
\bibitem{rubert} Kuratov Y., Arkhipov M. Adaptation of Deep Bidirectional Multilingual Transformers for Russian Language // arXiv. 2019. arXiv:1905.07213. URL: \url{https://arxiv.org/abs/1905.07213}.
\bibitem{wang} Wang S., Liu Y., Xu Y., Zhu C., Zeng M. Want To Reduce Labeling Cost? GPT-3 Can Help // Findings of ACL: EMNLP 2021. P. 4195--4205. DOI: 10.18653/v1/2021.findings-emnlp.354.
\bibitem{hinton} Hinton G., Vinyals O., Dean J. Distilling the Knowledge in a Neural Network // arXiv. 2015. arXiv:1503.02531. URL: \url{https://arxiv.org/abs/1503.02531}.
\bibitem{izmailov} Измайлов Л. С., Устинов А. М. Эмоциональный окрас комментариев на русском языке с помощью RuBERT моделей // Искусственный интеллект в промышленных, коммерческих, медицинских и финансовых приложениях. Вып. 1. М.: НИТУ «МИСИС», 2023. С. 3--8. EDN: GKTGPU.
\bibitem{beresnev} Береснев Д. В. Анализ подходов к использованию предобученных моделей в разработке корпоративных чат-ботов // Искусственный интеллект в промышленных, коммерческих, медицинских и финансовых приложениях. Вып. 1. М.: НИТУ «МИСИС», 2023. С. 68--77. EDN: OVKIHL.
\end{thebibliography}
\end{document}
